\section{Preliminaries}\label{sec:prelim}

\subsection{Memorization in the sense of Brown et al.}\label{sec:mem-def}
We use the framework of \citet[\S1]{Brown2021Memorization}. Fix a meta-distribution \(q\) over problem instances \(P\), each itself a distribution on labeled examples. A learning algorithm \(A\) takes a dataset \(X\) of \(n\) i.i.d.\ examples from \(P\) and outputs \(M=A(X)\). The (conditional) memorization is
\begin{equation}\label{eq:mem-def}
\Mem(A;q) \;:=\; I\!\left(X;\,M\,\middle|\,P\right),
\end{equation}
computed in the joint model \(P\sim q,\;X\sim P^{\otimes n},\;M=A(X)\). Equation~\eqref{eq:mem-def} measures how much the model reduces uncertainty about the realized training set beyond what is implied by knowing the data distribution.

\subsection{Federated release model}\label{sec:setup}
We extend \citet{Brown2021Memorization} to a partitioned, \emph{potentially non-IID} dataset. There are \(K\ge 2\) clients, and we let the meta-distribution \(q\) sample a \emph{joint problem instance}
\[
P\;=\;(P_1,\dots,P_K),
\]
where each \(P_k\) is a distribution over labeled examples. Conditional on \(P\), client \(k\) draws its local dataset \(X_k\sim P_k^{\otimes n_k}\), independently across \(k\); we write \(X=(X_1,\dots,X_K)\) and \(X_{-k}\) for everything except client~\(k\). The original single-distribution setting of \citet{Brown2021Memorization} is the IID-FL special case in which \(q\) is supported on the diagonal \(P_1=\dots=P_K\); when we refer below to ``conditioning on \(P\)'' we mean conditioning on the entire tuple. Cross-client heterogeneity is therefore part of the model from the outset, not an extension. One round of DP-FedAvg, parametrized by the clip bound \(C>0\) and the noise scale \(\sigma>0\), proceeds as follows.
Given the current iterate \(\theta\), each client computes a (possibly randomized) raw update
\[
V_k \;=\; V_k(X_k,\theta,R_k) \;\in\; \R^d,
\]
using fresh client randomness \(R_k\), and clips it to the ball of radius \(C\):
\begin{equation}\label{eq:clip}
U_k \;:=\; \mathrm{clip}_C(V_k) \;=\; V_k \cdot \min\!\Bigl\{1,\ \frac{C}{\|V_k\|_2}\Bigr\}, \qquad\text{so } \|U_k\|_2 \le C \text{ by construction.}
\end{equation}
The server averages the clipped updates and adds Gaussian noise:
\begin{equation}\label{eq:release}
    \noised \;:=\; \frac{1}{K}\sum_{k=1}^{K} U_k \;+\; Z, \qquad Z\sim\NN\!\left(0,\sigma^2 I_d\right),
\end{equation}
with \(Z\) drawn independently of all data and client randomness. The bound \(\|U_k\|_2\le C\) is part of the algorithm specification, not a hypothesis on \(V_k\); the analysis goes through for any choice of raw update map \(V_k\).
We make one assumption on the data-generating process:

\begin{assumption}[Conditional independence given the joint instance]\label{ass:indep}
Conditional on \(P=(P_1,\dots,P_K)\), the client datasets \(X_1,\dots,X_K\) are mutually independent, with \(X_k\sim P_k^{\otimes n_k}\) drawn from the \(k\)-th component. The client randomness \((R_1,\dots,R_K)\) is mutually independent and independent of \((X_1,\dots,X_K)\) and of \(P\). The server noise \(Z\) is drawn fresh, independent of \((X,R_1,\dots,R_K,P)\).
\end{assumption}

Assumption~\ref{ass:indep} generalizes the Brown et al.\ generative model from a single distribution to a tuple of per-client distributions, and is the independence model used implicitly in standard FL analyses (e.g., \citealt{Bornstein2022SWIFT}). The IID-FL special case (\(P_1=\dots=P_K\)) is recovered when \(q\) is concentrated on the diagonal; the non-IID case is the generic one and is what we have in mind throughout. The assumption restricts only the joint generation of client datasets given \(P\); it does \emph{not} restrict the per-client distributions \(P_k\) themselves, which may have arbitrarily different supports, label balance, or label-conditional structure across \(k\). All of the per-client memorization bounds below condition on the full tuple \(P\), so the inter-client heterogeneity is integrated into the conditioning rather than treated as adversarial perturbation.
